% T28 source: Примеры базовой модели.tex lines 258--397
\detailedsummaryheading{Full Parameters of the Illustrative Scenario}

We now move from abstract task sets to an illustrative scenario involving the
development of a small microservice. This is a synthetic scenario with
author-specified parameters, not an empirical study of a particular project.
The same project is considered in four variants to separate the influence of
two factors. In Variants 1--3, the task set remains unchanged while the
\emph{process configuration} changes: the parallelism and the agents' operating
mode. Variant 4 retains the best of these configurations but changes the
\emph{decomposition}: a large task on the critical path is split into subtasks.
This comparison makes it possible to assess separately the effect of work
organization and the effect of a finer project decomposition.

As before, the baseline effort of one unit of work is set to $X = 4$~hours.

\detailedexampleheading{Scenario Definition}

Consider a project to develop a small microservice for processing user
requests. The service accepts a request through an HTTP API, validates the
input, persists the state, initiates asynchronous processing, and provides
observability facilities. The project is compact enough to permit a complete
analysis, yet includes typical elements of production software development: an
API contract, data persistence, business rules, background processing, tests,
and infrastructure configuration.

For concreteness, assume that the service must support the following minimum
set of capabilities:
\begin{itemize}
  \item creating a request through the HTTP API;
  \item validating required fields and using a uniform error format;
  \item persisting the request and its status in a database;
  \item moving a request through the following states: created, accepted for
  processing, processed, and completed with an error;
  \item background processing with retries after transient errors;
  \item structured logs, metrics, a health check, and basic configuration;
  \item a suite of unit and integration tests;
  \item packaging and a minimal configuration for execution in a continuous
  integration and deployment environment.
\end{itemize}

For transparent parameterization of the illustrative scenario, we use notional
story points (SP) as units of task size. Let $2$~SP correspond to one unit of
$Z_i$. One story point then corresponds to $X/2 = 2$~hours, while the baseline
effort associated with one unit of $Z_i$ remains $X = 4$~hours. Subtasks with
an odd number of story points have fractional task sizes ($1$~SP $= 0.5\,Z$),
which is consistent with the condition $Z_i \in \mathbb{Q}^{+}$ in
Definition~\ref{def:params}. Calculations in Variants 1--3 use the aggregate tasks, whereas
Variant 4 directly uses the decomposition into subtasks.

The baseline decomposition of the project into aggregate tasks is shown in
Table~\ref{tab:practical-baseline}.

\begin{table}[ht!]
\centering
\caption{Baseline decomposition of the project into tasks}
\label{tab:practical-baseline}
\begin{tabularx}{\textwidth}{@{}c>{\raggedright\arraybackslash}Xcc@{}}
\toprule
$i$ & Task & Estimate, SP & Task size $Z_i$ \\
\midrule
1 & API contract and request validation & 4 & 2 \\
2 & Persistence layer (model, migrations, repository) & 6 & 3 \\
3 & Business logic (primary use cases and state transitions) & 8 & 4 \\
4 & Background worker (asynchronous processing and retries) & 6 & 3 \\
5 & Observability and configuration (logs, metrics, health check) & 4 & 2 \\
6 & Tests and integration checks & 6 & 3 \\
7 & Packaging and deployment configuration (Docker, CI) & 4 & 2 \\
\bottomrule
\end{tabularx}
\end{table}

The composition of the aggregate tasks is detailed in
Table~\ref{tab:practical-subtasks}. This detail is important for planning: at
the subtask level, it becomes apparent which work can be delegated to agents
independently and which work requires a contract to have been fixed in advance
or the result of another task. The detailed decomposition is used below to
identify dependencies introduced by graph coarsening and in Variant~4, where
the testing task is split into separate components.

\begin{table}[ht!]
\centering
\caption{Subtasks and story points}
\label{tab:practical-subtasks}
\begin{tabularx}{\textwidth}{@{}c>{\raggedright\arraybackslash}p{0.20\textwidth}>{\raggedright\arraybackslash}Xc@{}}
\toprule
$i$ & Task & Subtask & SP \\
\midrule
1 & API & Specify endpoints, DTOs, and the error format & 2 \\
1 & API & Implement data validation and error handling & 2 \\
2 & Persistence & Design database tables and migrations & 2 \\
2 & Persistence & Implement the repository and data-access methods & 3 \\
2 & Persistence & Add basic test data for local verification & 1 \\
3 & Business logic & Specify states and valid transitions & 2 \\
3 & Business logic & Implement the primary use cases & 4 \\
3 & Business logic & Handle domain-level errors & 2 \\
4 & Background worker & Implement retrieval of tasks for processing & 2 \\
4 & Background worker & Add retries and failure handling & 3 \\
4 & Background worker & Synchronize statuses with the primary model & 1 \\
5 & Observability & Add structured logs and metrics & 2 \\
5 & Observability & Configure settings and the health check & 2 \\
6 & Tests & Cover the business logic with unit tests & 2 \\
6 & Tests & Add integration tests for the API and persistence layer & 3 \\
6 & Tests & Verify background processing and retries & 1 \\
7 & Packaging & Prepare the Docker/service configuration & 2 \\
7 & Packaging & Add continuous-integration commands and a smoke test & 2 \\
\bottomrule
\end{tabularx}
\end{table}

The total project size in terms of the aggregate tasks is
\[
  \sum_{i=1}^{7} Z_i = 2 + 3 + 4 + 3 + 2 + 3 + 2 = 19.
\]
The equivalent total is $38$~SP, which, at $2$~hours per story point, also
yields $76$~hours of baseline work.

The no-AI completion time (one developer working sequentially) is
\[
  T_h = 19 \cdot 4 = 76~\text{h}.
\]

The tasks are linked by precedence relations. To plan parallel work, define a
graph $G=(V,E)$ whose vertices correspond to the aggregate tasks in
Table~\ref{tab:practical-baseline}. The dependencies are listed in
Table~\ref{tab:practical-dependencies}.

\begin{table}[ht!]
\centering
\caption{Dependency graph among the aggregate tasks}
\label{tab:practical-dependencies}
\begin{tabularx}{\textwidth}{@{}>{\raggedright\arraybackslash}p{0.29\textwidth}>{\raggedright\arraybackslash}p{0.12\textwidth}>{\raggedright\arraybackslash}X@{}}
\toprule
Task & Depends on & Work that can proceed in parallel \\
\midrule
1. API contract and request validation & --- & Initial task; establishes interfaces for the remaining work \\
2. Persistence layer & 1 & Can proceed in parallel with observability once the API is fixed \\
3. Business logic & 1, 2 & Can be partially designed after the API is fixed; implementation requires the persistence layer \\
4. Background worker & 2, 3 & Requires the data model and state-transition rules \\
5. Observability and configuration & 1 & Can proceed in parallel with the persistence layer and business logic \\
6. Tests and integration checks & 1, 2, 3, 4 & Unit tests can begin earlier; integration tests depend on the core components \\
7. Packaging and deployment & 1, 5 & Can be prepared in parallel once the interfaces and configuration are fixed \\
\bottomrule
\end{tabularx}
\end{table}

In edge notation, the graph is
\[
  E = \{1{\to}2,\; 1{\to}3,\; 1{\to}5,\; 1{\to}7,\; 2{\to}3,\; 2{\to}4,\; 3{\to}4,\; 3{\to}6,\; 4{\to}6,\; 5{\to}7\}.
\]

This graph implies a natural wave-based schedule:
\begin{enumerate}
  \item first, the API contract and basic DTOs are fixed;
  \item the persistence layer and observability can then be undertaken in
  parallel;
  \item once the persistence layer is complete, the full business logic becomes
  available, while packaging can be prepared in parallel;
  \item the background worker depends on the business logic and data model;
  \item integration tests and background-processing checks complete the
  critical path.
\end{enumerate}

The project thus combines parallelizable portions with strict dependencies. In
Variants 1--3, the aggregate values $Z_i$ and coefficients $k_i$ are used, but
completion time is determined not by balancing independent tasks, but by a
schedule on the graph in Table~\ref{tab:practical-dependencies}.

\paragraph{Actual Dependencies and Coarsening Artifacts.}
The coarsened graph conceals an important distinction. Some edges in
Table~\ref{tab:practical-dependencies} represent \emph{substantive}
dependencies: the business logic cannot rely on a data model that has not yet
been created, nor can the background worker rely on undefined state
transitions. Other edges arise solely from \emph{task coarsening}: they connect
entire tasks even though only individual subtasks are actually dependent. For
example, the dependence of Task 6 on Tasks 1--4 means that the test suite
contains checks for each component. Nevertheless, unit tests for the business
logic can begin immediately after Task 3, integration tests can begin once the
API and persistence layer are complete, and only the retry tests require a
completed background worker. Similarly, at the subtask level, the dependency
$3 \to 4$ pertains only to the 2-SP state-and-transition model within the 8-SP
task; the background worker does not require the remaining business-logic use
cases. A substantive dependency cannot be eliminated by rescheduling, whereas
a dependency induced by coarsening can be refined by replacing one vertex with
several subtasks. The dominance of the critical path in the calculations below
is therefore determined in part by the chosen graph granularity. Variant~4
quantifies this effect.

In Variants 1--3, the project and aggregate task set remain unchanged, but the
process configurations $\sigma = (P, r)$ differ: the number of parallel streams
$P$, the operating-mode assignment to tasks, and the stream assignment. Because the
coefficient $k_i = k_i^{(r)}$ depends on the ``task $\times$ tool $\times$
mode'' triple, interactive work in Variant 1, mixed-mode work in Variant 2, and
delegated work from detailed specifications in Variant 3 define different
vectors $\{k_i\}$ for the same task set. The tool and acceptance rules are fixed
within the scenario. Entire configurations are therefore compared, rather than
an isolated change in a single parameter. Variant~4 reproduces the
configuration of Variant~3 and changes only the decomposition, allowing its
effect to be assessed separately.

At level M4, competition among streams for developer attention is modeled
explicitly. Each task has a human component $h_i$ that includes task
specification, review, and the direction of corrections. A schedule is
considered feasible only if the human-attention intervals of different tasks do
not overlap and one agent remains continuously assigned to every task that has
started but has not yet been accepted. If a human-attention phase cannot begin
immediately, the agent remains blocked in the queue; other agents continue
working. In interactive mode, $h_i = k_i$ because the developer retains the
task context throughout its execution; in delegated mode, $h_i$ is
substantially smaller than $k_i$ but is nonzero because task specification and
acceptance are still required. The values of $h_i$, like those of $k_i$, are
assigned by expert judgment. Two simplifying assumptions are retained:
$C(P) \equiv 1$, meaning that cross-stream integration overhead is not modeled,
and $\gamma(P) \equiv 1$. The effect of attention switching is considered
separately when the variants are compared.
